\section{Python effect estimators: population, clipping and the last covariance step}
\label{exec20:pyfx-scope}
A bank evaluating retention offers must first decide whether it wants the
effect for all eligible customers, for customers historically offered the
intervention, or for a risk-defined population facing a future decision.
The same Python method name can hide different averaging, nuisance-fitting
and uncertainty choices. A reusable service must retain these choices in
its request and result. This section compares the already inspected EconML,
DoubleML, CausalML, Delicatessen and statsmodels paths, then develops two
conversion hazards: propensity clipping and the final covariance calculation.
These selected paths do not exhaust Python's effect-estimation ecosystem or
qualify continuous-dose, longitudinal or survival families by association.

\subsection{Five routes with different scientific responsibilities}
\label{exec20:pyfx-routes}
The selected EconML DML path residualizes outcome and treatment and fits an
effect model. Averaging predictions through an ATE interface retains that
model's restrictions; a misspecified constant projection does not become
the population effect merely because its output is called an average.
Its grouping for nuisance splits must also be distinguished from the
sampling law used for covariance. The earlier dynamic discussion shows why
unit histories, encounter order and bootstrap behavior need separate
inspection.~\citep{exec19_econml_source}

The inspected Python DoubleML IRM route makes the score, external nuisance
predictions, folds and repetitions central objects. Its ATE and ATTE choices
are distinct; the R implementation and other model families retain their
own contracts. Earlier matching of seven configurations with fixed identical
nuisance predictions and folds is useful bounded runtime evidence, but does
not qualify new learners, unrestricted dependence or a complete package
distribution. In particular, nuisance vectors must be joined by stable row
keys when another repository prepares them.~\citep{runtime_doubleml_scope}

The selected CausalML T-learner instead fits arm-specific outcomes and reports
contrasts in a defined treatment-label order. An average over prediction rows
depends on the chosen prediction population. Its row bootstrap is not a
borrower-history or common-date bootstrap, and prediction percentiles and
the default ATE interval formula are different objects. An initial binary
treatment profile avoids silently extending that limited review to every
multiarm variance calculation.~\citep{exec14_causalml_primary}

Delicatessen is attractive when the proposal needs transparent stacked
estimating equations and joint nuisance/target inference. The inspected
implementation differentiates the supplied equations and builds a sandwich
from their contribution columns. Those columns encode the sampling units;
the routine does not reconstruct household or calendar dependence from
labels supplied elsewhere. Its allowed pseudoinverse can make a calculation
available without establishing that the target is locally identified.
Analytic small-model derivatives, rank checks and a defensible sampling law
are therefore part of the proposed use.~\citep{v13_sw_delicatessen,v131_sw_delicatessen}

The selected statsmodels \texttt{TreatmentEffect} path supplies five
point-treatment procedures: inverse probability weighting (IPW), regression
adjustment (RA), augmented IPW (AIPW), AIPW with weighted outcome fits, and
IPW regression adjustment. Its documented outcome model is linear OLS.
It combines target, outcome and selection moments for inference where
applicable. This is a useful parametric comparison route, with a different
nuisance-fitting architecture from a cross-fitted flexible learner.
The inspected revision is \texttt{8278e2d218cc} of 6 October 2026; the
following claims concern its selected source, not an unverified later
release.~\citep{v13_sw_statsmodels,exec20_python_effect_source}

For a proposal, the choice is consequently conditional. Use a declared DML
or IRM profile when its score and effect model match the question; retain
arm-specific fits for a separately qualified CATE comparison; use explicit
stacked equations or the parametric statsmodels paths as interpretable
baselines. Reimplementing all five would duplicate model, variance and
failure semantics before a measured need exists. Original engines behind
typed boundaries are the preferred candidates, with executable admission
limited to profiles that later pass their source, environment and behavior
checks. DoWhy, zEpid, causallib and CausalPy remain complementary leads for
target-specific inspection; naming them here grants no capability finding.

\subsection{All customers, treated customers and untreated customers}
\label{exec20:pyfx-population}
Let $A\in\{0,1\}$ denote the historical offer, $X$ baseline information,
$Y(a)$ a contribution outcome under intervention $a$, and
$\tau(X)=\mathbb E[Y(1)-Y(0)\mid X]$. Under consistency, conditional
exchangeability and positivity in the chosen eligible population,
$\tau(X)$ is identified by the corresponding conditional outcome contrast.
Three averages still differ:
\begin{equation}
 \theta_{\rm all}=\mathbb E[\tau(X)],\quad
 \theta_{\rm treated}=\frac{\mathbb E[p(X)\tau(X)]}{\mathbb E[p(X)]},\quad
 \theta_{\rm untreated}=\frac{\mathbb E[(1-p(X))\tau(X)]}{\mathbb E[1-p(X)]},
 \label{exec20:pyfx-targets}
\end{equation}
where $p(X)=\Pr(A=1\mid X)$. The latter two expressions follow by conditioning
on $X$ and using the historical assignment probabilities. They refer to
historical populations; an effect for a future targeting rule requires that
rule's population and information set.

For an exact authored example, suppose low- and high-risk customers are
equally common. Their offer probabilities are $1/4$ and $3/4$, and their
conditional effects are USD 20 and USD 80. Then
$\theta_{\rm all}=50$, $\theta_{\rm treated}=65$ and
$\theta_{\rm untreated}=35$. At a stipulated incremental cost of USD 55,
the corresponding net averages are USD $-5$, USD 10 and USD $-20$.
Selecting the treated-population answer for an all-customer proposal could
therefore reverse the investment story. These are imposed-law calculations,
not observed benefits or a reason to use historical treatment as a future
policy.~\citep{exec20_python_effect_reference}

The inspected statsmodels methods expose these target distinctions through
\texttt{effect\_group}. RA averages fitted potential outcomes over the
selected rows. IPW-RA changes both arm-fitting weights and prediction rows:
for the treated target, for example, treated-arm weights are one and
control-arm weights are propensity odds. AIPW and AIPW-WLS share
target-specific potential-outcome terms. Their joint moments reorder
selection scores to match the grouped outcome rows. An adapter must preserve
these masks and row identities through every nuisance component, rather
than rely on position after independent sorting. Arbitrary user-weight
arrays are outside the initial proposed profile until their point and
inference paths are separately traced.~\citep{exec20_python_effect_source}

\subsection{Clipping an estimated probability changes the score}
\label{exec20:pyfx-clipping}
Consider the ordinary all-population AIPW score with proposed outcome
functions $m_a(X)$ and a working propensity $q(X)$:
\begin{equation}
 \phi_q=m_1(X)-m_0(X)+\frac{A\{Y-m_1(X)\}}{q(X)}
       -\frac{(1-A)\{Y-m_0(X)\}}{1-q(X)}.
 \label{exec20:pyfx-score}
\end{equation}
Write $\mu_a(X)=\mathbb E[Y\mid A=a,X]$, and suppose $0<q(X)<1$.
Taking the expectation conditional on $X$ gives the exact identity
\begin{equation}
 \mathbb E[\phi_q\mid X]-\{\mu_1(X)-\mu_0(X)\}
 =(p-q)\left\{\frac{\mu_1-m_1}{q}
                 +\frac{\mu_0-m_0}{1-q}\right\}.
 \label{exec20:pyfx-bias}
\end{equation}
The right side vanishes if $q=p$, or if both outcome functions equal their
conditional means. Replacing a correct $p$ by a clipped $q$ generally removes
the first route to a zero bias. This is a score identity; estimation error,
sample splitting and asymptotic remainder control need their own argument.

At the inspected source revision, \texttt{ps\_bounds} defaults to
$(0.001,0.999)$ and probabilities are clipped in both point calculations
and selected joint-GMM moment paths. AIPW-WLS also uses probability-derived
outcome weights; the final target calculation still uses the same
population-specific augmented terms. Clipping can reduce numerical
extremes, but it neither establishes positivity in the source population
nor manufactures the rare untreated observations that inference needs.
\citep{exec20_python_effect_source}

To isolate the consequence, take a constant true propensity
$p=1999/2000$, a clipped value $q=999/1000$, conditional means
$\mu_0=20,\mu_1=80$, and deliberately incorrect working outcome functions
$m_0=m_1=0$. The unclipped expected score is 60. The clipped expectation
is $69970/999$, approximately 70.04, with bias $10030/999$.
At a stipulated USD 65 cost, the true mean margin is USD $-5$, whereas
the clipped-score margin is $5035/999$, approximately USD 5.04. Correct
outcome functions would restore the expected score to 60 even with this
clipping. No native package was run, and no claim is made that a fitted
OLS model in this example would select the imposed zero functions.
The example shows what a blanket statement about robustness after clipping
would have to rule out.~\citep{exec20_python_effect_reference}

A proposed service should therefore return the original probability range,
clipping bounds and counts, target weights and overlap diagnostics. If it
trims observations or changes to an overlap-weighted target, that is another
population contract. A sensitivity display across clipping choices is useful
only when the underlying score and target are named; a visually stable
coefficient is not a proof of identification or calibrated uncertainty.

\subsection{Follow inference through to the final moment covariance}
\label{exec20:pyfx-covariance}
The statsmodels class explicitly states that its preliminary outcome
covariance option does not control the final GMM result. The selected five
wrappers call joint GMM without an alternative \texttt{weights\_method};
the default is \texttt{cov}. The final covariance code reads stored fit
options and recomputes moment covariance. HAC is a separate path requiring
a lag choice and kernel; the selected wrapper signatures do not expose that
choice merely by accepting a preliminary robust covariance label.
Thus an HC1 label on an intermediate weighted regression cannot establish
household clustering or common-date uncertainty for the final causal effect.
\citep{exec20_python_effect_source,v131_sw_statsmodels}

Joint nuisance moments can account for estimation of the specified
parametric nuisance model under their sampling and regularity conditions.
They do not justify silently substituting arbitrary machine-learned
predictions, a dependent panel law or a different target score. Record the
full moment vector, derivative matrix, numerical residuals and rank. An
optimizer returning a result establishes a numerical event, not the required
orthogonality, model correctness or sampling approximation.

One small source detail illustrates why that numerical record matters.
The selected covariance routine centers a moment matrix by
\texttt{moms.mean()}, a scalar mean over all entries. Let each row
$m_i\in\mathbb R^k$ be one observation's moments, let
$\bar m=n^{-1}\sum_i m_i$, and set $c=k^{-1}\mathbf1^{\mathsf T}\bar m$.
Denote the second-moment matrices after column and scalar centering by
$C_{\rm col}$ and $C_{\rm scalar}$. Expansion gives
\begin{equation}
 C_{\rm scalar}-C_{\rm col}
   =(\bar m-c\mathbf1)(\bar m-c\mathbf1)^{\mathsf T}.
 \label{exec20:pyfx-centering}
\end{equation}
Cross terms vanish because the column-centered rows sum to zero.
If all moment means are exactly zero, the two constructions agree.
At nonzero residual means they can differ, especially when moment scaling
and an ill-conditioned derivative amplify a numerical discrepancy.
This identity is a qualification question, not a measured software defect
or an allegation of invalid intervals in every fit.
\citep{exec20_python_effect_source}

For the authored moment rows $(1,0)$ and $(3,0)$,
$\bar m=(2,0)$ and $c=1$. The two matrices are
\begin{equation}
 C_{\rm col}=\begin{pmatrix}1&0\\0&0\end{pmatrix},\qquad
 C_{\rm scalar}=\begin{pmatrix}2&-1\\-1&1\end{pmatrix}.
 \label{exec20:pyfx-centering-example}
\end{equation}
The difference is the outer product of $(1,-1)$ with itself. These rows
deliberately do not solve zero-mean moment equations. A conformance test
should therefore include both an exact root and controlled nonzero
residuals, recording the derivative and covariance response rather than
concluding from this example that a particular native estimate is wrong.
\citep{exec20_python_effect_reference}

\subsection{A public prototype that tests meaning before scale}
\label{exec20:pyfx-design}
The first comparison should freeze one binary intervention, units, an
eligible population and independent sampling units. Use the exact target
and clipping laws above, followed by generated heterogeneous outcomes with
known assignment, intentional nuisance misspecification and weak overlap.
Fix row keys, folds, repetitions, predictions and treatment-level order.
For each engine, compare the target expression and retained population
before comparing numbers. Evaluate uncertainty only under its declared law;
then add separately justified clustered or temporal regimes.

The same contract should travel through library, CLI and MCP requests:
target population, treatment contrast, nuisance identities, weighting and
uncertainty settings enter together. The response preserves point estimate,
uncertainty, diagnostics and unavailable states separately. Under an
incorrect model or unsupported dependence law, a successful tool call
must not become a qualified causal claim. Package upgrades reopen the
relevant conversion checks, including row sorting, clipping defaults,
target labels, covariance construction and error handling.

The proposed comparison can measure whether this semantic discipline helps
analysts detect consequential mistakes before an expensive bank experiment.
Setup and dependency maintenance, nuisance training, memory, conversions,
diagnostics and repeated inference are cost categories to measure, not
current prices or savings. Complete distribution provenance and rights,
target-specific native fixtures and broader original-source coverage remain
open. Five reviewed implementations supply concrete choices and tests;
they do not establish a completed ecosystem survey or universal inference.
